% =====================================================================
\section{Thuật toán Randomized SVD}\label{sec:algorithm}
% =====================================================================

\subsection{Ý tưởng cốt lõi}\label{subsec:rsvd-idea}

Với ma trận $A \in \mathbb{R}^{m\times n}$ kích thước cực lớn, ta chỉ thực sự
cần thông tin về \emph{không gian con chính} $\operatorname{range}(A_k)$ --
nơi tập trung gần như toàn bộ năng lượng. Nếu xác định được một ma trận
trực giao $Q \in \mathbb{R}^{m\times \ell}$ ($\ell$ nhỏ, $\ell \ll \min(m,n)$)
sao cho
\begin{equation}\label{eq:Q-good}
    A \;\approx\; Q\, Q^T A,
\end{equation}
thì việc tính SVD của $A$ quy về tính SVD của ma trận
$B = Q^T A \in \mathbb{R}^{\ell\times n}$ -- kích thước nhỏ hơn rất nhiều.

\emph{Quan sát then chốt} của Halko, Martinsson và Tropp~\cite{halko2011}:
nếu lấy $\Omega \in \mathbb{R}^{n\times \ell}$ là ma trận ngẫu nhiên với các
phần tử i.i.d.\ $\mathcal{N}(0, 1)$, thì hầu như chắc chắn (với xác suất rất
gần $1$) các cột của $Y = A\Omega$ phủ kín không gian con chính của $A$. Cụ
thể, với $\ell = k + p$ trong đó $p \ge 5$ là tham số \emph{over-sampling},
ta thu được xấp xỉ \eqref{eq:Q-good} với sai số kiểm soát được dưới dạng kỳ
vọng (xem \cref{subsec:error-bound}).

\subsection{Mô tả thuật toán}\label{subsec:rsvd-pseudocode}

Thuật toán được trình bày chính thức trong \cref{alg:rsvd}. Đầu vào là ma trận
$A$, hạng đích $k$, tham số over-sampling $p$ (mặc định $p = 10$) và số bước
power iteration $q$ (mặc định $q = 1$--$2$).

\begin{algorithm}[ht]
\caption{Randomized SVD (Halko--Martinsson--Tropp 2011)}\label{alg:rsvd}
\begin{algorithmic}[1]
\Require Ma trận $A \in \mathbb{R}^{m\times n}$; hạng đích $k$;
         over-sampling $p$; số bước power iteration $q$; seed ngẫu nhiên.
\Ensure  $U \in \mathbb{R}^{m\times k}$,
         $s \in \mathbb{R}^{k}$ (giảm dần),
         $V^T \in \mathbb{R}^{k\times n}$, sao cho $A \approx U \operatorname{diag}(s) V^T$.
\State $\ell \gets \min(k + p,\; n)$
    \Comment{kích thước sketch}
\State Sinh $\Omega \in \mathbb{R}^{n\times \ell}$ với các phần tử
       i.i.d.\ $\mathcal{N}(0,1)$
    \Comment{ma trận chiếu Gauss}
\State $Y \gets A\,\Omega$
    \Comment{$Y \in \mathbb{R}^{m\times \ell}$: ``sketch'' không gian cột}
\For{$i = 1$ to $q$} \Comment{power iteration: tăng độ tách phổ}
    \State $Q,\, R \gets \text{QR}(Y)$
        \Comment{tái trực giao hóa chống suy thoái số}
    \State $Z \gets A^T Q$
    \State $Q_2,\, R_2 \gets \text{QR}(Z)$
    \State $Y \gets A\, Q_2$
\EndFor
\State $Q,\, R \gets \text{QR}(Y)$
    \Comment{cơ sở trực chuẩn cho $\operatorname{range}(Y)$}
\State $B \gets Q^T A$
    \Comment{$B \in \mathbb{R}^{\ell\times n}$: ``thu nhỏ'' ma trận gốc}
\State $\widetilde{U},\, s,\, V^T \gets \text{SVD}(B)$
    \Comment{SVD đầy đủ trên ma trận nhỏ}
\State $U \gets Q\, \widetilde{U}$
    \Comment{lift về không gian gốc}
\State \Return $U[:, :k]$, $s[:k]$, $V^T[:k, :]$
    \Comment{truncate về hạng $k$}
\end{algorithmic}
\end{algorithm}

\subsection{Cơ sở toán học của mỗi bước}\label{subsec:rsvd-justify}

\paragraph{Tại sao $Y = A\Omega$ chứa không gian con chính của $A$?}
Viết SVD $A = U\Sigma V^T$, ta có
\[
    A \Omega \;=\; U\Sigma V^T \Omega \;=\; U\Sigma\, \widetilde{\Omega},
\]
với $\widetilde{\Omega} = V^T \Omega \in \mathbb{R}^{n\times \ell}$ vẫn có
phân bố Gauss chuẩn (do $V^T$ là trực giao). Mỗi cột của $Y$ là tổ hợp tuyến
tính ngẫu nhiên của các cột của $U\Sigma$, với trọng số tỉ lệ với $\sigma_i$.
Vì $\sigma_1 \ge \dots \ge \sigma_k \gg \sigma_{k+1}$ (giả định phổ có gap),
các thành phần $u_1, \dots, u_k$ thống trị trong $Y$. Khi $\ell \ge k$, với
xác suất gần $1$, $\operatorname{span}(Y) \approx \operatorname{span}(u_1, \dots, u_k)$.

\paragraph{Vai trò của over-sampling $p$.}
Nếu chỉ lấy $\ell = k$, ma trận ngẫu nhiên $\widetilde{\Omega}$ có thể có hạng
không đầy đủ trên không gian con đầu (xác suất thấp nhưng khác $0$). Thêm
$p$ cột ngẫu nhiên \emph{ngoài} $k$ giúp triệt tiêu xác suất này: với $p = 10$,
xác suất thất bại xuống dưới $10^{-15}$~\cite{halko2011}.

\paragraph{Vai trò của power iteration.}
Nếu phổ kỳ dị của $A$ suy giảm chậm
(ratio $\sigma_{k+1}/\sigma_k$ gần $1$), thuật toán cơ sở (không power
iteration) có thể bắt được thành phần $\sigma_{k+1}, \sigma_{k+2}, \dots$ một
cách nhầm lẫn. Quan sát: ma trận $(A A^T)^q A$ có cùng véc-tơ kỳ dị với $A$
nhưng phổ bị \emph{nén lũy thừa} -- giá trị kỳ dị thứ $i$ trở thành
$\sigma_i^{2q+1}$, do đó tỉ số trở thành $(\sigma_{k+1}/\sigma_k)^{2q+1}$. Mỗi
bước power iteration tăng độ tách phổ theo \emph{cấp số nhân}. Trong cài đặt
thực tế, ta áp dụng QR sau mỗi bước nhân (dòng 5--8 của \cref{alg:rsvd}) để
tránh tràn số do mũ cao.

\paragraph{Tại sao $B = Q^T A$ giữ nguyên thông tin kỳ dị?}
Nếu $\operatorname{range}(Q) \approx \operatorname{range}(A_k)$ thì
$Q Q^T A \approx A_k$, tức $Q^T A$ chứa toàn bộ thông tin về $\Sigma$ và $V$
của $A_k$. Tính SVD $B = \widetilde{U} \Sigma V^T$ cho ta trực tiếp $\Sigma, V$;
còn $U = Q\widetilde{U}$ là $\widetilde{U}$ ``nâng lên'' không gian
$\mathbb{R}^m$ ban đầu thông qua $Q$.

\subsection{Phân tích độ phức tạp}\label{subsec:complexity}

Đặt $\ell = k + p$. Chi phí từng bước của \cref{alg:rsvd}:

\begin{table}[ht]
\centering
\caption{Chi phí của các bước trong \cref{alg:rsvd}, với
$A$ dày so với $A$ thưa có $\mathrm{nnz}$ phần tử khác $0$.}
\label{tab:complexity}
\begin{tabular}{lll}
\toprule
\textbf{Bước} & \textbf{$A$ dày} & \textbf{$A$ thưa} \\
\midrule
Sinh $\Omega$              & $\mathcal{O}(n\ell)$
                           & $\mathcal{O}(n\ell)$ \\
$Y = A\,\Omega$            & $\mathcal{O}(mn\ell)$
                           & $\mathcal{O}(\mathrm{nnz}\cdot\ell)$ \\
$q$ bước power iteration   & $\mathcal{O}(q\cdot mn\ell)$
                           & $\mathcal{O}(q\cdot \mathrm{nnz}\cdot\ell)$ \\
QR của $Y$                 & $\mathcal{O}(m\ell^2)$
                           & $\mathcal{O}(m\ell^2)$ \\
$B = Q^T A$                & $\mathcal{O}(mn\ell)$
                           & $\mathcal{O}(\mathrm{nnz}\cdot\ell)$ \\
SVD của $B$ (kích thước $\ell\times n$) & $\mathcal{O}(n\ell^2)$
                                        & $\mathcal{O}(n\ell^2)$ \\
\bottomrule
\end{tabular}
\end{table}

Tổng chi phí (bỏ qua hằng số):
\begin{equation}\label{eq:rsvd-cost}
\boxed{\;\mathcal{C}_{\mathrm{RSVD}} \;=\;
    \mathcal{O}\bigl((mn + (m + n)\ell)\,\ell\bigr)
\;\text{(dày)},
\quad
    \mathcal{O}\bigl(\mathrm{nnz}\cdot\ell + (m+n)\ell^2\bigr)
\;\text{(thưa)}.}
\end{equation}

So với SVD cổ điển $\mathcal{O}(mn\min(m,n))$: khi $\ell = k + p \ll \min(m, n)$,
tỷ số gia tốc lý thuyết là $\Theta(\min(m,n)/\ell)$. Với $m = n = 10^5$ và
$k = 500$, tỉ lệ này lên tới $\sim 200\times$. Trên ma trận thưa, lợi thế còn
lớn hơn nhiều do tận dụng được cấu trúc $\mathrm{nnz} \ll mn$.

\subsection{Sai số xấp xỉ}\label{subsec:error-bound}

\begin{theorem}[Halko--Martinsson--Tropp~\cite{halko2011}]\label{thm:hmt-bound}
Cho $A \in \mathbb{R}^{m\times n}$ và $\Omega \in \mathbb{R}^{n\times (k+p)}$
là ma trận ngẫu nhiên Gauss chuẩn, $p \ge 2$. Gọi $Q$ là ma trận trực giao
xây dựng từ QR của $A\Omega$. Khi đó kỳ vọng sai số xấp xỉ thỏa mãn
\begin{equation}\label{eq:hmt-bound}
    \mathbb{E}\bigl[\, \|A - Q Q^T A\|_F\,\bigr]
       \;\le\; \biggl(1 + \frac{k}{p - 1}\biggr)^{1/2}\,
               \biggl(\sum_{i > k}\sigma_i^2\biggr)^{1/2}.
\end{equation}
\end{theorem}

\begin{remark}[So với Eckart--Young]
Theo \cref{thm:eckart-young}, xấp xỉ hạng-$k$ tốt nhất có sai số chính bằng
$\sqrt{\sum_{i>k}\sigma_i^2}$. \cref{thm:hmt-bound} cho thấy Randomized SVD
chỉ \emph{lệch một hằng số nhỏ} so với optimum. Với $p = 10$ và $k = 100$,
hệ số trước căn là $\sqrt{1 + 100/9} \approx 3.5$. Khi thêm $q$ bước power
iteration, bound thực tế còn chặt hơn đáng kể (xem định lý 10.7
của~\cite{halko2011}).
\end{remark}

\subsection{Tiêu chuẩn chọn $k$ trong ngữ cảnh Randomized SVD}\label{subsec:k-selection}

Trong ngữ cảnh Randomized SVD, ta chỉ tính được $\ell = k_{\max} + p$ giá trị
kỳ dị đầu tiên. Để chọn $k$ thỏa mãn \eqref{eq:energy-criterion} cần biết
\emph{tổng} năng lượng $\|A\|_F^2$. Tuy nhiên, đại lượng này có thể tính
\emph{trực tiếp} từ $A$ mà không cần SVD đầy đủ:

\begin{equation}\label{eq:total-energy-direct}
    \|A\|_F^2 \;=\; \sum_{i, j} a_{ij}^2.
\end{equation}

Với ma trận thưa lưu dưới định dạng CSR/CSC, công thức
\eqref{eq:total-energy-direct} chỉ tốn $\mathcal{O}(\mathrm{nnz}(A))$ -- một
phép cộng các bình phương của các phần tử khác $0$. Đây là một điểm tinh tế
\emph{bị bỏ qua trong nhiều cài đặt}: nếu lấy
$\sum_{i \le k_{\max}}\sigma_i^2$ làm tổng năng lượng (như notebook gốc ban
đầu của nhóm), ta sẽ overestimate tỉ lệ năng lượng giữ lại vì phần đuôi
$i > k_{\max}$ bị bỏ qua hoàn toàn.

\medskip
\noindent\textbf{Quy trình chọn $k$ (được dùng trong toàn bộ thực nghiệm):}
\begin{enumerate}[label=(\arabic*)]
    \item Chọn ngưỡng $\eta \in (0, 1)$;
    \item Tính $E = \|A\|_F^2$ trực tiếp từ \eqref{eq:total-energy-direct};
    \item Chạy Randomized SVD bậc $k_{\max}$ đủ lớn, thu
          $\sigma_1 \ge \sigma_2 \ge \dots \ge \sigma_{k_{\max}}$;
    \item Tính $S_k = \sum_{i=1}^{k}\sigma_i^2$. Chọn
          $k^\star = \min\{k\colon S_k/E \ge \eta\}$;
    \item Nếu $S_{k_{\max}}/E < \eta$, tăng $k_{\max}$ và lặp lại.
\end{enumerate}

\subsection{Điều kiện break-even cho nén dữ liệu}\label{subsec:break-even}

Lưu trữ biểu diễn nén $(U_k, \Sigma_k, V_k)$ chiếm $k(m + n + 1)$ phần tử
\texttt{float64}. So với ma trận gốc:

\begin{itemize}
    \item \emph{Ma trận dày}: nguyên bản chiếm $mn$ phần tử, do đó cần
    $k(m+n+1) < mn$, hay
    \begin{equation}\label{eq:be-dense}
        k \;<\; k_{be}^{\text{dày}} \;=\; \frac{mn}{m + n + 1}
        \;\approx\; \frac{\min(m, n)}{2}.
    \end{equation}
    Điều kiện này dễ dàng thỏa mãn trong thực tế.

    \item \emph{Ma trận thưa} (lưu CSR): nguyên bản chiếm
    $\sim \mathrm{nnz}$ phần tử (cộng index overhead), do đó
    \begin{equation}\label{eq:be-sparse}
        k \;<\; k_{be}^{\text{thưa}} \;=\; \frac{\mathrm{nnz}(A)}{m + n + 1}.
    \end{equation}
    Với ma trận \emph{rất thưa} (mật độ thấp), $k_{be}^{\text{thưa}}$ rất nhỏ
    và việc nén bằng SVD hiếm khi có lợi về dung lượng.
\end{itemize}

\eqref{eq:be-sparse} là kết quả thực tế quan trọng được nhóm phát hiện trong
quá trình thực nghiệm với MovieLens 25M (\cref{sec:experiments}): mặc dù
$k^\star$ đạt năng lượng cao, $k^\star \gg k_{be}^{\text{thưa}}$ nên SVD
\emph{không nén được dữ liệu thưa} -- biểu diễn nén còn lớn hơn ma trận gốc.

\subsection{Tóm tắt cài đặt}\label{subsec:impl-summary}

Nhóm cài đặt \cref{alg:rsvd} trong tệp \texttt{src/rsvd.py} dùng thuần
\texttt{NumPy}/\texttt{SciPy}. Hàm chính có chữ ký:

\begin{lstlisting}[caption={Chữ ký hàm Randomized SVD tự cài
(\texttt{src/rsvd.py}).},label={lst:rsvd-api}]
def randomized_svd(A, k, p=10, q=2, random_state=None):
    """
    Truncated SVD via random projection (Halko-Martinsson-Tropp).

    Returns
    -------
    U  : ndarray, shape (m, k)
    s  : ndarray, shape (k,)
    Vt : ndarray, shape (k, n)
    """
\end{lstlisting}

Cài đặt hỗ trợ \emph{cả} ma trận dày (\texttt{numpy.ndarray}) lẫn ma trận thưa
(\texttt{scipy.sparse}): phép nhân $A\Omega$ và $A^T Q$ được thực hiện thông
qua hàm \texttt{\_matmat}/\texttt{\_rmatmat} bao bọc, gọi đúng routine cho
từng loại đầu vào. Toàn bộ mã chỉ phụ thuộc \texttt{numpy} và \texttt{scipy};
không có lời gọi nào tới SVD đầy đủ trên $A$.
